% ============================================================================
%  math_version/sections/part3a_utility_core.tex  --  Part III, Section 7
%  "The acquisition objective" (utility core).
%  \input-ready fragment: no preamble/documentclass. Uses the shared macros
%  from preamble.tex. Domain-neutral (pure math / ML) rendering; symbols and
%  mathematics identical to the source. The deep proofs (conditioning,
%  divergence, predictive-mean equivalence) live in the sibling Part III-deep
%  section and are only cited here.
% ============================================================================

\section{The acquisition objective}\label{sec:utility}

The active-learning loop proceeds sequentially: after each count observation the
model is refit, and the next input to query is the one whose observation is
expected to teach us the most about the latent function. ``Teach us the most'' is
made precise as an \emph{information gain} --- an expected reduction in the entropy
of what we do not yet know about the latent function and its rate. This section
builds that utility on the generative model and variational posterior of
Parts~I--II: the Poisson-lognormal predictive of the count, its entropy, and the
two utilities considered here.

Throughout, the candidate (query) input is $x^{*}$; a count is the random variable
$R$ with realizations $r$ (this is the count written $\spk$ elsewhere in the
document --- we keep it as $r$ here to match the count-domain notation $\rmax$ and
the Laplace mode $\gmode$). The mean rate at input $x$ is
$\rate(x)=\exp(\gain\,\lat(x)+\latz)$ with latent $\lat(x)$, gain $\gain$, and
offset $\latz$; the count is $R\mid\rate(x)\sim\Poi(\rate(x))$. It is convenient to
work with the \emph{log-rate} $\lgr(x)=\gain\,\lat(x)+\latz$, so $\rate=e^{\lgr}$,
whose GP posterior at any input is Gaussian with the log-rate moments
%
\begin{equation}
  \mug \;=\; \gain\,\pmean(x^{*})+\latz,
  \qquad
  \varg \;=\; \gain^{2}\,\pvar(x^{*}),
  \label{eq:logratemoments}
\end{equation}
%
the affine push-forward of the latent posterior moments $\pmean(x^{*})$,
$\pvar(x^{*})$ read off the variational posterior~\eqref{eq:predmoments}.

\begin{remark}[Two utilities, one idea]\label{rem:two-utilities}
We consider \emph{two} acquisition functions. Both express the same idea ---
\emph{utility$(x^{*})=$ information the observation $R$ at $x^{*}$ gives about the
latent function $=$ a drop in entropy} --- and they differ only in \emph{which}
uncertainty they score (Table~\ref{tab:two-utilities}). The \textbf{standard
utility} $\Ustd$ scores the rate \emph{at $x^{*}$ itself}; the
\textbf{distribution-aware utility} $\Uda$ scores the rates at the inputs
$x\sim\px$ we ultimately want to predict. The standard utility depends only on the
marginal posterior at $x^{*}$; the distribution-aware utility couples $x^{*}$ to the
input distribution $\px$ through cross-covariance. $\Uda$ is the richer object, but
it carries a norm-exploitation divergence (\S\ref{subsec:landscape},
Thm.~\ref{thm:divergence}) that a finite candidate pool masks.
\end{remark}

\begin{table}[t]
  \centering
  \small
  \begin{tabular}{@{}lll@{}}
    \toprule
      & \textbf{Standard} $\Ustd$ & \textbf{Distribution-aware} $\Uda$\\
    \midrule
    Formula      & $\Hmarg(x^{*})-\E[\Hnoise(x^{*})]$
                 & $\Hmarg(x^{*})-\E_{x\sim\px}[\Hcond(x^{*})]$\\
    As MI        & $I\!\big(\rate(x^{*});R\mid x^{*},\mathcal D\big)$
                 & $I\!\big(\rate(X);R\mid X,x^{*},\mathcal D\big),\ X\sim\px$\\
    Scores       & rate \emph{at $x^{*}$ itself}
                 & rates at \emph{inputs} $x\sim\px$\\
    Needs        & marginal moments at $x^{*}$
                 & marginal moments \emph{and} cross-cov.\ to $\px$\\
    \bottomrule
  \end{tabular}
  \caption{The two utilities. Both are ``information gain $=$ entropy reduction'',
  about different targets.}
  \label{tab:two-utilities}
\end{table}

% ---------------------------------------------------------------------------
\subsection{Utility as expected entropy reduction and mutual information}
\label{subsec:mi}

Fix a candidate $x^{*}$ and the data $\mathcal D=\{(x_i,r_i)\}$. The utility is the
mutual information between the count $R$ we would observe at $x^{*}$ and the latent
function we want to learn. Precisely what we want to learn about is the difference
between the two utilities.

\paragraph{The distribution-aware objective.}
Let $Z:=(X,\rate(X))$ with $X\sim\px$ an input drawn from the input distribution
and $\rate(X)$ its rate. Choosing $x^{*}$ to maximize the information $R$ carries
about $Z$,
%
\begin{equation}
  \Uda(x^{*}) \;:=\; I\!\big(Z;R\mid x^{*},\mathcal D\big)
             \;=\; I\!\big((X,\rate(X));\,R\mid x^{*},\mathcal D\big).
  \label{eq:da-objective}
\end{equation}
%
The chain rule for mutual information splits this into a term for \emph{which} input
we imagine and a term for its \emph{rate}:
%
\begin{equation}
  I\!\big((X,\rate(X));R\mid x^{*},\mathcal D\big)
  \;=\;
  \underbrace{I(X;R\mid x^{*},\mathcal D)}_{=\,0}
  \;+\;
  I\!\big(\rate(X);R\mid X,x^{*},\mathcal D\big).
  \label{eq:mi-chain}
\end{equation}
%
The first term vanishes: \emph{which} input $X$ we intend to ask about is
independent of the count $R$ observed at $x^{*}$, given $x^{*}$ and $\mathcal D$.
Hence
%
\begin{equation}
  \Uda(x^{*})
  \;=\; I\!\big(\rate(X);R\mid X,x^{*},\mathcal D\big)
  \;=\; \E_{x\sim\px}\!\big[\,I\!\big(\rate(x);R\mid x,x^{*},\mathcal D\big)\big].
  \label{eq:da-expectation}
\end{equation}

\paragraph{Entropy decomposition.}
For a fixed input $x$, the per-input mutual information decomposes into a
marginal-entropy term minus an expected conditional-entropy term,
%
\begin{equation}
  I\!\big(\rate(x);R\mid x,x^{*},\mathcal D\big)
  \;=\;
  H(R\mid x^{*},\mathcal D)
  \;-\;
  \E_{\rate(x)}\!\big[\,H(R\mid x^{*},\rate(x),\mathcal D)\big],
  \label{eq:mi-entropy-split}
\end{equation}
%
using (i) $H(R\mid x,x^{*},\mathcal D)=H(R\mid x^{*},\mathcal D)$ --- the marginal
of $R$ at $x^{*}$ does not depend on which input we intend to predict --- and
(ii) that conditioning on the latent value $\rate(x)$ leaves the expected entropy
$\E_{\rate(x)}[H(R\mid x^{*},\rate(x),\mathcal D)]$. Averaging
\eqref{eq:mi-entropy-split} over $x\sim\px$ and naming the two pieces gives the
distribution-aware utility as a difference of entropies,
%
\begin{equation}
  \Uda(x^{*}) \;=\; \Hmarg(x^{*}) \;-\; \E_{x\sim\px}\!\big[\Hcond(x^{*})\big],
  \label{eq:Uda}
\end{equation}
%
with
%
\begin{align}
  \Hmarg(x^{*}) &:= H(R\mid x^{*},\mathcal D), \label{eq:def-Hmarg-cond}\\
  \Hcond(x^{*}) &:= \E_{x\sim\px,\ \lat(x)\sim q}\!\big[\,H(R\mid x^{*},\lat(x),\mathcal D)\big].
  \label{eq:def-Hcond}
\end{align}
%
By the symmetry of mutual information, $\Uda(x^{*})$ equals the reduction in
uncertainty about the rates $\{\rate(x):x\sim\px\}$ produced by observing $R$ at
$x^{*}$ --- the precise sense of ``how much does querying $x^{*}$ teach us about the
latent function over the input distribution.''

\begin{remark}[Two expectations inside $\Hcond$]\label{rem:double-exp}
$\Hcond$ in~\eqref{eq:def-Hcond} carries \emph{two} expectations: over inputs
$x\sim\px$ and over the posterior latent $\lat(x)\sim q$. The latent is integrated
because $\rate(x)=\exp(\gain\lat(x)+\latz)$ is itself uncertain --- only the
posterior of $\lat(x)$ is known, not its value. Replacing the inner expectation by
an evaluation at the mean $\lat(x)=\pmean(x)$ is \emph{biased}; the necessity of the
inner sampling is the subject of \S\ref{subsec:why-sample}.
\end{remark}

The \emph{standard} utility instead scores the rate at the query point itself. It
is the special case in which the ``input we predict'' is $x^{*}$, so $\Uda$'s outer
expectation collapses and $\Hcond$ becomes the aleatoric Poisson noise entropy at
$x^{*}$; we develop it in \S\ref{subsec:standard} once $\Hmarg$ is in hand.

% ---------------------------------------------------------------------------
\subsection{The marginal count entropy \texorpdfstring{$\Hmarg$}{H\_marg}}
\label{subsec:hmarg}

Both utilities share the marginal predictive entropy of the count at $x^{*}$,
%
\begin{equation}
  \Hmarg(x^{*})
  \;=\; -\sum_{r=0}^{\infty} p(r\mid x^{*},\mathcal D)\,\log p(r\mid x^{*},\mathcal D),
  \label{eq:Hmarg}
\end{equation}
%
where the marginal predictive integrates the Poisson likelihood against the
Gaussian posterior of the log-rate $\lgr(x^{*})\sim\Gauss(\mug,\varg)$
(cf.~\eqref{eq:logratemoments}):
%
\begin{equation}
  p(r\mid x^{*},\mathcal D)
  \;=\; \int \Poi(r\mid e^{\lgr})\,\Gauss(\lgr\mid\mug,\varg)\,d\lgr
  \;=\; \frac{1}{r!}\int \exp\!\big(r\lgr - e^{\lgr}\big)\,\Gauss(\lgr\mid\mug,\varg)\,d\lgr .
  \label{eq:marg-predictive}
\end{equation}
%
This integral has no closed form. It is evaluated by a \textbf{Laplace expansion} of
the integrand about its mode $\gmode$ (the mode of the log-integrand
$r\lgr - e^{\lgr} - (\lgr-\mug)^{2}/2\varg$):
%
\begin{equation}
  \log p(r\mid x^{*},\mathcal D)
  \;\approx\;
  \gmode\,r - e^{\gmode}
  - \frac{(\gmode-\mug)^{2}}{2\varg}
  - \tfrac12\log\!\big(1+\varg\,e^{\gmode}\big)
  - \log r! .
  \label{eq:laplace-logp}
\end{equation}

\begin{lemma}[Closed-form Laplace mode via Lambert~$W$]\label{lem:mode}
The stationarity condition $r-e^{\lgr}-(\lgr-\mug)/\varg=0$ rearranges to
$\lgr=r\varg+\mug-\varg e^{\lgr}$. Substituting $w=\varg e^{\lgr}$ gives
$w+\log w=\log\varg+r\varg+\mug$, i.e.\ $w=\LW\!\big(\varg\,e^{\,r\varg+\mug}\big)$
with $\LW$ the principal branch of the Lambert~$W$ function. Hence the mode is
closed-form,
%
\begin{equation}
  \gmode \;=\; r\varg+\mug-\LW\!\big(\varg\,e^{\,r\varg+\mug}\big).
  \label{eq:laplace-mode}
\end{equation}
\end{lemma}

The infinite sum in~\eqref{eq:Hmarg} decays fast (rates are low) and is truncated at
$\rmax$ (\S\ref{subsec:landscape} handles the truncation artifact). At very small
variance the Laplace form is unstable, and one falls back to the exact Poisson
log-pmf $\log p(r)=r\mug-e^{\mug}-\log r!$ (used for $\varg<10^{-6}$). The mode is
computed in log-space --- $\LW(e^{y})$ with $y=\log\varg+r\varg+\mug$,
Newton-iterated with the derivative $dW/dy=W/(1+W)$ --- so no $e^{r\varg+\mug}$ is
ever formed (overflow-safe and differentiable in $x^{*}$).

% ---------------------------------------------------------------------------
\subsection{The standard utility}
\label{subsec:standard}

The standard acquisition is the mutual information between the count and the latent
rate \emph{at the query point itself},
%
\begin{equation}
  \Ustd(x^{*})
  \;=\; \Hmarg(x^{*}) - \E[\Hnoise(x^{*})]
  \;=\; H(R\mid x^{*},\mathcal D) - \E_{\rate\sim\text{post.}}\!\big[H(R\mid\rate)\big]
  \;=\; I\!\big(\rate(x^{*});R\mid x^{*},\mathcal D\big).
  \label{eq:Ustd}
\end{equation}
%
It scores ``how much would a count at $x^{*}$ pin down the rate at $x^{*}$'' --- it
depends on the marginal GP moments at $x^{*}$ only, with no $\px$ and no
cross-covariance. It is exactly $\Uda$ with the predicted input taken to be the
query point, so the second term is the \emph{aleatoric} Poisson noise entropy rather
than an epistemic conditional entropy.

\subsubsection{The aleatoric noise entropy \texorpdfstring{$\E[\Hnoise]$}{E[H\_noise]} in closed form}
\label{subsubsec:hnoise}

Given a rate $\rate$, the entropy of $\Poi(\rate)$ is
$\Hnoise(\rate)=\rate\,(1-\log\rate)+\E_{r\sim\Poi(\rate)}[\log r!]$. Writing
$\lgr=\log\rate\sim\Gauss(\mug,\varg)$ and using the lognormal moment identities
%
\begin{equation}
  \E[e^{\lgr}] = \exp\!\big(\mug+\tfrac12\varg\big) = \fbar(x^{*}),
  \qquad
  \E[\lgr\,e^{\lgr}] = \exp\!\big(\mug+\tfrac12\varg\big)\,(\mug+\varg),
  \label{eq:lognormal-identities}
\end{equation}
%
(the first is the expected rate $\fbar$ of~\eqref{eq:fbar}), the first piece
integrates in closed form:
%
\begin{equation}
  \E_{\lgr}\!\big[\rate(1-\log\rate)\big]
  = \E_{\lgr}\!\big[e^{\lgr}(1-\lgr)\big]
  = \exp\!\big(\mug+\tfrac12\varg\big)\big(1-(\mug+\varg)\big)
  = -\exp\!\big(\mug+\tfrac12\varg\big)(\mug+\varg-1).
  \label{eq:hnoise-firstpiece}
\end{equation}
%
The remaining $\E[\log r!]$ is taken under the marginal predictive
$p(r\mid x^{*},\mathcal D)$ of~\eqref{eq:marg-predictive}. Hence
%
\begin{equation}
  \E[\Hnoise(x^{*})]
  \;=\; -\exp\!\big(\mug+\tfrac12\varg\big)(\mug+\varg-1)
        \;+\; \sum_{r} p(r\mid x^{*},\mathcal D)\,\log r! .
  \label{eq:Hnoise}
\end{equation}
%
This is the closed form for $\E[\Hnoise]$.
Assembling $\Ustd$ then reads the marginal latent moments
$\pmean(x^{*}),\pvar(x^{*})$, transforms them to $\mug,\varg$
by~\eqref{eq:logratemoments}, and returns $\Ustd=\Hmarg-\E[\Hnoise]$.

% ---------------------------------------------------------------------------
\subsection{The distribution-aware utility}
\label{subsec:da}

The distribution-aware utility~\eqref{eq:Uda} reduces uncertainty about the rates at
the inputs we actually want to predict. It weights a candidate $x^{*}$ by how
informative its observation is about the rates at typical inputs --- i.e.\ by the
cross-covariance between $x^{*}$ and inputs $x\sim\px$ (``distribution-aware'' $=$
aware of the input distribution $\px$). Evaluating $\Hcond$ needs the predictive
moments of $\lat(x^{*})$ \emph{after a hypothetical observation} $\lat(x)=\lat_i$ at
an input $x$.

\subsubsection{Conditional GP moments via Gaussian conditioning}
\label{subsubsec:condmoments}

Let $\pv(x)=\Ktil^{-1}\kvec(x)$ be the SVGP projection and $q(\latt)=\Gauss(\vm,\vV)$
the variational posterior on the inducing values. The marginal latent posterior
moments at $x$ and $x^{*}$ are the projected moments of~\eqref{eq:predmoments},
$\pmean(x)=\pv(x)\transpose\vm$, $\pvar(x)=s(x)+\pv(x)\transpose\vV\pv(x)$
(and likewise at $x^{*}$), where $s(x)=\kvec(x,x)-\pv(x)\transpose\kvec(x)$ is
the prior Schur complement. The object that couples the two points is the
\textbf{prior conditional covariance} of $\lat(x),\lat(x^{*})$ given the inducing
values:
%
\begin{equation}
  \ccond \;=\; \kvec(x,x^{*}) - \pv(x)\transpose\Ktil\,\pv(x^{*}) .
  \label{eq:util-ccond}
\end{equation}
%
By the law of total covariance the full posterior cross-covariance is
%
\begin{equation}
  \Sigma_{x^{*}} \;=\; \ccond + \pv(x)\transpose\vV\,\pv(x^{*})
             \;=\; \kvec(x,x^{*}) + \pv(x)\transpose(\vV-\Ktil)\,\pv(x^{*}) ,
  \label{eq:util-crosscov}
\end{equation}
%
and standard Gaussian conditioning on the joint of
$(\lat(x),\lat(x^{*}))$ gives the conditional latent moments at $x^{*}$ after
observing $\lat(x)=\lat_i$,
%
\begin{align}
  \mu_{\mathrm{cond}}(x^{*})     &= \pmean(x^{*}) + \frac{\Sigma_{x^{*}}}{\pvar(x)}\,\big(\lat_i-\pmean(x)\big),
    \label{eq:util-mucond}\\[2pt]
  \sigma^{2}_{\mathrm{cond}}(x^{*}) &= \pvar(x^{*}) - \frac{\Sigma_{x^{*}}^{2}}{\pvar(x)} .
    \label{eq:util-sigcond}
\end{align}
%
These are transformed to log-rate space, $\gain\,\mu_{\mathrm{cond}}+\latz$ and
$\gain^{2}\sigma^{2}_{\mathrm{cond}}$, and fed to the same entropy
functional~\eqref{eq:Hmarg} to give the per-input conditional entropy
$H(R\mid x^{*},\lat(x),\mathcal D)$. The full Gaussian-conditioning derivation
(augmented inducing matrix, block inverse, and the equivalence to the rank-1
posterior update) is the deep-layer result~\eqref{eq:condmoments}; and the
conditional mean~\eqref{eq:util-mucond} obtained by conditioning on the joint
posterior coincides with the mean obtained by augmenting the inducing set with the
hypothetical observation and re-integrating (Thm.~\ref{thm:predmean-equiv}). Neither
is re-derived here.

\begin{remark}[The cross-covariance pitfall: $\ccond$ must be included]\label{rem:cond-correct}
Forming the joint over $[x;x^{*}]$ and reading the \emph{full} posterior covariance
between the two points already yields $\ccond+\pv\transpose\vV\pv^{*}=\Sigma_{x^{*}}$
--- the $\ccond$ term of~\eqref{eq:util-ccond} included. The pitfall is to instead
assemble the $\pv\transpose\vV\pv^{*}$-only cross-covariance by hand, dropping
$\ccond$: this is biased outside the inducing hull. Taking
$\mu_{\mathrm{cond}},\sigma^{2}_{\mathrm{cond}}$ from the blocks of the complete
joint covariance (with a floor $\sigma^{2}_{\mathrm{cond}}\ge10^{-8}$) avoids it.
\end{remark}

\paragraph{Monte-Carlo estimator (one $\lat$ per input).}
Given pre-drawn samples $\{x_i\}_{i=1}^{N_{\mathrm{mc}}}\sim\px$:
%
\begin{enumerate}\itemsep2pt
  \item $\Hmarg(x^{*})=$ entropy~\eqref{eq:Hmarg} at the marginal moments
        $\mug,\varg$ of $x^{*}$.
  \item For $i=1,\dots,N_{\mathrm{mc}}$: read $(\pmean(x_i),\pvar(x_i))$; draw
        $\lat_i=\pmean(x_i)+\sqrt{\pvar(x_i)}\,\varepsilon,\ \varepsilon\sim\Gauss(0,1)$
        (sampled variant; the deterministic variant sets $\lat_i=\pmean(x_i)$); form
        the conditional moments \eqref{eq:util-mucond}--\eqref{eq:util-sigcond} at
        $x^{*}$; set $\Hcond^{(i)}=$ entropy at $(\gain\mu_{\mathrm{cond}}+\latz,\ \gain^{2}\sigma^{2}_{\mathrm{cond}})$.
  \item $\Hcond=\tfrac{1}{N_{\mathrm{mc}}}\sum_i\Hcond^{(i)}$;\quad
        $\Uda=\Hmarg-\Hcond$.
\end{enumerate}
%
The input posteriors and the reparameterized draws $\lat_i$ do not depend on
$x^{*}$, so they are held fixed (detached from the gradient) when differentiating in
$x^{*}$, avoiding $O(N_{\mathrm{mc}})$ graph growth; the conditional moments at
$x^{*}$ remain differentiable in $x^{*}$. The two variants: the deterministic one
takes $\lat_i=\pmean(x_i)$ (a reproducible sanity check ---
\S\ref{subsubsec:det}); the sampled one is the correct unbiased estimator
(\S\ref{subsec:why-sample}).

\subsubsection{The deterministic special case}
\label{subsubsec:det}

Setting $\lat_i=\pmean(x)$ (the posterior mean) makes the innovation
$\lat_i-\pmean(x)=0$, so~\eqref{eq:util-mucond}--\eqref{eq:util-sigcond} collapse
to $\mu_{\mathrm{cond}}=\pmean(x^{*})$ (mean unchanged) and
%
\begin{equation}
  \sigma^{2}_{\mathrm{cond}}(x^{*}) \;=\; \pvar(x^{*})\,(1-\corr^{2}),
  \qquad
  \corr^{2} \;=\; \frac{\Sigma_{x^{*}}^{2}}{\pvar(x)\,\pvar(x^{*})},
  \label{eq:varreduce}
\end{equation}
%
where $\corr$ is the posterior latent correlation between $\lat(x)$ and
$\lat(x^{*})$. In log-rate coordinates the conditioning map is
$(\mug,\varg)\mapsto(\mug,\varg(1-\corr^{2}))$: same mean, variance scaled by
$1-\corr^{2}$. Writing $\Hmarg(\cdot,\cdot)$ for the count
entropy~\eqref{eq:Hmarg} as a function of the log-rate moments, the utility is the
entropy difference
%
\begin{equation}
  \Uda(x^{*}) \;=\; \Hmarg(\mug,\varg) \;-\; \Hmarg\!\big(\mug,\ \varg(1-\corr^{2})\big).
  \label{eq:da-deterministic}
\end{equation}
%
So the deterministic $\Uda$ depends on exactly three quantities: the operating
point $\mug$, the marginal log-rate variance $\varg$ (how much entropy to start
with), and $\corr^{2}$ (the fraction of that variance conditioning removes).
$\corr^{2}$ is governed by the angle between $x^{*}$ and the conditioning input in
$\Cmat$-space: a small angle gives $\corr^{2}\!\to\!1$ (near-total variance
removal), a large angle $\corr^{2}\!\to\!0$ (no conditioning). This same $\corr^{2}$
sets the sampling bias of \S\ref{subsec:why-sample}.

% ---------------------------------------------------------------------------
\subsection{Why we sample the latent instead of using its mean}
\label{subsec:why-sample}

$\Hcond$ is a double expectation (Remark~\ref{rem:double-exp}), over $x\sim\px$ and
over $\lat(x)\sim q$. The tempting shortcut replaces the inner expectation by
evaluation at the mean,
%
\begin{equation}
  \widetilde{H}_{\mathrm{cond}}(x^{*})
  \;=\; \E_{x\sim\px}\!\big[\,H\big(R\mid x^{*},\lat(x)=\pmean(x),\mathcal D\big)\big].
  \label{eq:mean-shortcut}
\end{equation}

\begin{proposition}[The mean shortcut is biased]\label{prop:bias}
$\E_{\lat(x)}[H(R\mid x^{*},\lat(x),\mathcal D)]\neq H(R\mid x^{*},\lat(x)=\pmean(x),\mathcal D)$
in general. The predictive mean of $\lat(x^{*})$ is linear in the conditioning
value, $\E[\lat(x^{*})\mid\lat(x),\mathcal D]=\pmean(x^{*})+\alpha\,(\lat(x)-\pmean(x))$
with regression coefficient $\alpha=\mathrm{Cov}(\lat(x^{*}),\lat(x)\mid\mathcal D)/\pvar(x)$,
while the predictive \emph{variance} does not depend on the value of $\lat(x)$.
Entropy is nonlinear in the predictive mean, so a second-order Taylor expansion of
$H$ about $\lat(x)=\pmean(x)$, with $\E[\lat(x)-\pmean(x)]=0$, leaves the leading
bias
%
\begin{equation}
  \E_{\lat(x)}[H] - H\big|_{\lat(x)=\pmean(x)}
  \;\approx\; \tfrac12\,\alpha^{2}\,\pvar(x)\;\frac{\partial^{2}\Hmarg}{\partial\pmean^{2}}\Big|_{x^{*}} .
  \label{eq:bias-raw}
\end{equation}
%
Since $\mathrm{Cov}(\lat(x^{*}),\lat(x))=\corr\,\sqrt{\pvar(x^{*})\pvar(x)}$, we
have $\alpha^{2}\,\pvar(x)=\corr^{2}\,\pvar(x^{*})$, and the $\pvar(x)$ cancels:
%
\begin{equation}
  \mathrm{Bias} \;\approx\; \frac{\corr^{2}}{2}\;\pvar(x^{*})\;\frac{\partial^{2}\Hmarg}{\partial\pmean^{2}}\Big|_{x^{*}} .
  \label{eq:bias}
\end{equation}
\end{proposition}

The bias vanishes \emph{only} when $\corr=0$ ($x$ uninformative about $x^{*}$ ---
but then the utility contribution is $\approx0$ anyway), or $\pvar(x^{*})=0$
(nothing left to learn), or $\partial^{2}\Hmarg/\partial\pmean^{2}=0$ (entropy
linear in the mean --- false for the Poisson-lognormal). Notably $\pvar(x)=0$ is
\emph{not} a zero-bias condition: it cancels in~\eqref{eq:bias}.

\begin{remark}[The bias is largest exactly where it matters]\label{rem:bias-matters}
Bias $\propto\corr^{2}\,\pvar(x^{*})$ is maximal when the sampled input $x$ is most
correlated with $x^{*}$ (most relevant to the utility) \emph{and} when $x^{*}$ is
most uncertain (where the acquisition should be steering us) --- i.e.\ largest as
$\corr\to1$. A systematic bias does not shrink with more Monte-Carlo samples: it
converges to the \emph{wrong} value and corrupts the optimizer's direction.
Sampling $\lat$ is unbiased and its variance falls as $O(1/N)$. Hence: sample
$\lat$.
\end{remark}

\paragraph{One $\lat$ per input is optimal (nested MC).}
With $g(x,\lat)=H(R\mid x^{*},\lat,\mathcal D)$, let $\tau^{2}=\mathrm{Var}_x[\E_{\lat}g]$
(between-input) and $\E_x[\sigma^{2}_x]$ (within-input). For a fixed budget of $M$
entropy evaluations, using $N_{\lat}$ latent samples per input gives variance
$(N_{\lat}\tau^{2}+\E_x[\sigma^{2}_x])/M$, strictly worse for $N_{\lat}>1$.
Spending the budget on \emph{more inputs}, one $\lat$ each, is optimal.

\caveat{The magnitude of this bias is a provisional estimate. Pushed to log-rate
space the curvature carries a gain factor,
$\partial^{2}\Hmarg/\partial\pmean^{2}=\gain^{2}\,\partial^{2}\Hmarg/\partial\lgr^{2}$,
so $\mathrm{Bias}\propto\gain^{2}$; for a small gain (e.g.\ $\gain\approx0.03$) the
bias is far below the MC standard error, and the deterministic and sampled $\Uda$
nearly coincide. The regime guide $\gain\ll1$ negligible / $\gain\sim1$ moderate /
$\gain\gg1$ essential, and the specific numbers, are plausibility arguments, not
settled results. The core claim --- bias~\eqref{eq:bias} does not vanish with $N$
--- is solid; the magnitude estimate is provisional. The \emph{correct} estimator
remains the sampled one regardless.}

% ---------------------------------------------------------------------------
\subsection{The entropy landscape over input space}
\label{subsec:landscape}

\paragraph{Monotone shape --- utility is not pure epistemic uncertainty.}
The count entropy $\Hmarg(\mug,\varg)$ increases monotonically in \emph{both}
arguments: with $\mug$ (higher predicted rate) and with $\varg$ (more
GP/epistemic uncertainty). A candidate therefore scores high for a \emph{mixture}
of reasons --- it is predicted to have a high rate \emph{and} it is epistemically
uncertain. The utility is thus not a pure epistemic-uncertainty acquisition;
predicted rate is baked in. This is a deliberate modelling choice: a normalized
arc-cosine kernel would force $\Kf(x,x)=1$, removing the norm channel so that
utility tracked only angular alignment (pure epistemic), but the held-out
predictive correlation then drops by $\sim$25\% (a fitted example, $\Mind=100$:
$0.79\!\to\!0.59$) --- input norm is genuinely informative, so the kernel keeps the
norm dependence.

\paragraph{The $\rmax$ truncation artifact.}
With a fixed $\rmax=100$ the Laplace sum~\eqref{eq:Hmarg} misses the probability
mass above $\rmax$ once the Poisson peak climbs past it, and $\Hmarg$
\emph{collapses to $\approx0$} for $\mug\gtrsim\log100\approx4.6$ (and at lower
$\mug$ for large $\varg$). This is a numerical artifact, not real entropy: the
computable region is roughly triangular in $(\mug,\varg)$. A cheap pre-check
without running the Laplace sum is $z_{\mathrm{safe}}=(\log\rmax-\mug)/\sqrt{\varg}$
--- $z_{\mathrm{safe}}>2$ safe, $<1.5$ unreliable, $<0.5$ broken. Inputs from the
pool sit at $z_{\mathrm{safe}}\gg10$, always safe; trouble arises only for
artificially amplified inputs. An adaptive choice of $\rmax$, sized from the moments
so the truncation always covers the Poisson mass, removes the collapse.

\paragraph{Norm-exploitation divergence.}
Because the arc-cosine kernel is positively $1$-homogeneous
($\Kf(\scal x,y)=\scal\,\Kf(x,y)$ when the bias $\sigz=0$), scaling
$x^{*}\to\scal x^{*}$ scales $\pmean(x^{*})\sim\scal$ and $\pvar(x^{*})\sim\scal^{2}$,
so $\Hmarg$ grows while a directionally aligned conditioning input keeps
$\corr\approx1$ and $\sigma^{2}_{\mathrm{cond}}\approx0$. Under free-space gradient
ascent $\Uda$ then grows $\sim\scal^{1.9}$ (numerically) --- the optimizer
``diverges'' by amplifying input norm rather than finding angularly informative
inputs. This is intrinsic to the homogeneous kernel, \emph{not} a bug; the full
statement and the $\sigz^{2}>0$ damping (alignment peak strictly at $\scal=1$) are
Thm.~\ref{thm:divergence} in the deep layer. Over a finite \emph{pool} of candidate
inputs (discrete selection) norms are bounded and this never triggers.

\caveat{\textbf{Utility accuracy ceiling (known limitation).} At high rate the two
utilities fail in opposite directions through the same $\rmax=100$ truncation:
$\Ustd$ can run to $+\infty$ while $\Uda$ silently collapses toward $0$. A rate
guard ($f_{\max}$, default $100$) keeps the operating point in the accurate regime;
a dedicated look-up table for the $\Uda$ path was not built. This is a limitation of
the fixed-$\rmax$ path, cured on the standard path by the adaptive choice of
$\rmax$.}
